\textbf{Plano dimensional de obra nueva.} La Figura~\ref{fig:plano-recinto-proyectado} fija un módulo técnico nuevo de 10,0$\times$7,0 m interiores y 3,0 m de altura libre, adyacente a administración, de acceso técnico independiente. Es una condición de obra y distribución, no una superficie atribuida al armario existente. Las zonas separan cómputo, comunicaciones, UPS, baterías, climatización, trabajo y respaldo; generación y combustible permanecen en el exterior. El plano ejecutivo comprobará implantación, suelo, evacuación, incendio, refrigerante y distancias reglamentarias antes de ejecutar las obras \parencite[RT-06.01/03/06, pp.~14--15]{pucv-transversal}.

\begin{figure}[htbp]
\centering
\setstretch{1}
\begin{tikzpicture}[x=1.35cm,y=1.35cm,
  zona/.style={draw=SynNavyLight,fill=SynSurface,thick},
  texto/.style={font=\fontsize{9.5}{11.5}\selectfont,align=center},
  cota/.style={{Latex}-{Latex},thin,draw=SynNavyLight}]
  \draw[zona] (0,0) rectangle (10,7);
  \draw[zona] (0,0) rectangle (2.5,3.8);
  \draw[zona] (0,3.8) rectangle (2.5,7);
  \draw[zona] (7.5,0) rectangle (10,3.8);
  \draw[zona] (7.5,3.8) rectangle (10,7);
  \draw[zona] (2.5,1.5) rectangle (7.5,7);
  \draw[dashed,draw=SynNavyLight] (5,1.5)--(5,4.1);
  \draw[dashed,draw=SynNavyLight] (2.5,5)--(7.5,5);
  \node[texto,text width=29mm] at (1.25,5.4) {Baterías A\\2 cadenas\\40 bloques/cadena\\Aire separado de TI};
  \node[texto,text width=29mm] at (8.75,5.4) {Baterías B\\2 cadenas\\40 bloques/cadena\\Aire separado de TI};
  \node[texto,text width=28mm] at (1.25,1.9) {UPS A\\93T60\\Distribución A};
  \node[texto,text width=28mm] at (8.75,1.9) {UPS B\\93T60\\Distribución B};
  \node[texto,text width=28mm] at (3.75,2.9) {Cómputo\\24 U\\8 U ocupadas};
  \node[texto,text width=28mm] at (6.25,2.9) {Comunicaciones\\12 U\\10 U ocupadas};
  \draw[draw=SynNavyLight,fill=white,thick] (3.4,5.95) rectangle (4.6,6.65);
  \draw[draw=SynNavyLight,fill=white,thick] (5.4,5.95) rectangle (6.6,6.65);
  \node[texto] at (4,6.3) {Clima A};
  \node[texto] at (6,6.3) {Clima B};
  \draw[-{Latex},thick,draw=SynBlue] (4,5.95)--(4,4.2);
  \draw[-{Latex},thick,draw=SynBlue] (6,5.95)--(6,4.2);
  \draw[-{Latex},thick,dashed,draw=SynNavyLight] (3,4.2) |- (3.4,6.3);
  \draw[-{Latex},thick,dashed,draw=SynNavyLight] (7,4.2) |- (6.6,6.3);
  \node[texto,fill=SynSurface,text width=60mm] at (5,4.75) {Recinto común: 27,5 m$^2$\\TI y unidades interiores; sin muro entre zonas};
  \node[texto,fill=SynSurface,text width=55mm] at (5,3.9) {Suministro continuo / retorno discontinuo};
  \node[texto,text width=61mm] at (5,.75) {Circulación, acceso independiente\\frente útil $\geq1,20$ m};
  \draw[cota] (0,7.3)--(10,7.3);
  \node[texto,fill=white,inner sep=1pt] at (5,7.3) {10,0 m interiores};
  \draw[cota] (-.3,0)--(-.3,7);
  \node[texto,rotate=90,fill=white,inner sep=1pt] at (-.3,3.5) {7,0 m interiores};
  \node[texto,draw=SynNavyLight,fill=SynSurface,text width=65mm,anchor=north west,inner sep=2mm] at (0,-.45)
    {\textbf{Exterior A/B segregado}\\Cada grupo: zona $\geq6,0\times4,0$ m.\\Depósito propio: 1.500 L nominales.\\Volumen útil $\geq1.190$ L; contención neta $\geq1.650$ L.};
  \node[texto,draw=SynNavyLight,fill=SynSurface,text width=55mm,anchor=north east,inner sep=2mm] at (10,-.45)
    {\textbf{Trabajo y respaldo separados}\\Área de trabajo externa $\geq3,0\times2,0$ m.\\Stock local apagado en espacio $\geq2,0\times1,5$ m.\\Bancada y ensayos fuera de sala.};
\end{tikzpicture}
\caption{Plano de distribución proyectada del recinto técnico y zonas asociadas}
\label{fig:plano-recinto-proyectado}
\FuenteFigura{Elaboración de Synaptix: cotas y condiciones de obra nueva según \parencite[RT-06.03, p.~14]{pucv-transversal} y área de instalación según \parencite[secc.~2.2.6, p.~21]{p9-mext-manual-2024}. Las divisiones discontinuas delimitan usos dentro del recinto común; las flechas representan circulación de aire. El plano ejecutivo desarrolla holguras, seguridad de refrigerante y tratamiento de los recintos energéticos separados.}
\end{figure}


La superficie se obtiene de dos zonas energéticas de $2(2,5\times7)=35$ m$^2$, un recinto común de $17,5+10=27,5$ m$^2$ para cómputo/comunicaciones y ambas unidades interiores de climatización, y 7,5 m$^2$ de circulación frontal. Las divisiones de uso dentro del recinto común no son muros estancos; mantienen comunicación de aire permanente. La reserva de superficie responde a acceso, separación y mantenimiento de cuatro gabinetes de baterías y dos caminos; no se convierte en una expansión del parque TI ni en capacidad cloud local. La implantación puede compactarse al verificar equipos y holguras, conservando soporte y acceso. Esta adaptación mantiene el núcleo principal en nube y justifica el espacio de continuidad del sitio secundario \parencites[secc.~6.1, p.~14; RT-15.01, p.~27]{pucv-transversal}[cap.~15, RT-06.01]{caso07}.

El módulo técnico completo tiene 70 m$^2$ y 210 m$^3$, con recintos separados; el compartimento común TI/climatización ocupa 27,5 m$^2$ y 82,5 m$^3$; cada lateral energético reserva 2,5$\times$7,0 m, dividido en UPS de 2,5$\times$2,2 m y bancos de 2,5$\times$4,8 m. Cada gabinete de cadena se gira a huella de 0,8$\times$2,0 m y se dispone consecutivamente en el lateral: coordenadas interiores A $(0,2;2,5)$--$(1,0;4,5)$ y $(0,2;4,7)$--$(1,0;6,7)$; B es su reflejo respecto de $x=5$. El frente largo de servicio queda hacia el pasillo de 1,30 m; la variante debe admitir servicio frontal y los 0,20 m traseros reservados o se redimensionará. La puerta de cada banco, de 1,20 m libres, abre hacia fuera en la fachada norte, al extremo del pasillo, sin barrer equipos ni bloquear salida. Cada cadena de cuarenta bloques ocupa un gabinete de proyecto de hasta 2,0$\times$0,8$\times$1,8 m, con cuatro niveles de diez bloques, protección y acceso propio; se instalan dos gabinetes por camino. Su masa máxima de diseño es $1.756+300=2.056$ kg, donde 300 kg es cota propia para estructura, barras y accesorios. Con base continua de 1,60 m$^2$, la presión media es $2.056(9,80665)/1,60=12,6015$ kN/m$^2$. Se prescribe solera calculada para al menos 25 kN/m$^2$ en toda zona energética, verificando además cargas concentradas, traslado, anclaje sísmico y capacidad del suelo; una media admisible no sustituye el cálculo de apoyos. No se apilan gabinetes ni se instala la reserva chilena en esa huella \parencite[p.~1]{csbHRL12540WRA260131}.

El área de instalación se verifica sobre ese compartimento común, de $5,0\times5,5=27,5$ m$^2$, con las dos unidades interiores físicamente dentro y suministro/retorno como muestra la figura. El manual s-MEXT G00 de junio de 2024 exige al menos 21 m$^2$ para 022 F2; no se suman los recintos separados de bancos, UPS o circulación para acreditar ese mínimo. Se prescribe acceso técnico categoría C, sin ocupación permanente y con un máximo de dos personas durante intervención: $2/27,5=0,0727<0,10$ personas/m$^2$. El frente de servicio propio de 1,20 m supera los 850 mm del manual, conservando además sus holguras de circulación según versión OVER/UNDER. La selección final verificará carga R32 por circuito y presente en el recinto, longitud de tubería, ubicación, circulación y requisitos aplicables a ambas unidades; el cotejo de superficie no acredita por sí solo seguridad integral del refrigerante ni capacidad sensible/latente \parencite[secc.~2.2.5--2.2.7, p.~21]{p9-mext-manual-2024}.

Cada UPS reserva 0,50$\times$0,82$\times$1,75 m y 235,5 kg de equipo sin bancos, según su ficha. Cada gabinete AR3104 de cómputo/comunicaciones tiene huella real $0,600\times1,070$ m y altura 1,1985 m, con 24 U por gabinete. La elevación RIO-4 conserva 8 U/10 U reservadas; la solera TI se calcula para al menos 10 kN/m$^2$ y reacciones concentradas. Se ubican en sus zonas de $2,5\times3,5$ m, frente a $2,5\times(1,070+1,20+1,20)=8,675$ m$^2$ de envolvente de gabinete y accesos anterior/posterior; quedan 0,03 m de holgura longitudinal en esta envolvente conservadora. Respecto del origen interior de la figura, cómputo ocupa $(3,45;2,70)$--$(4,05;3,77)$ y comunicaciones $(5,95;2,70)$--$(6,55;3,77)$; los frentes miran al acceso sur. Los accesos anteriores abarcan $y=1,50$--$2,70$ y los posteriores $y=3,77$--$4,97$, sin ocupar los climatizadores en $y=5,95$--$6,65$. El acceso lateral se conserva sin compartirlo con el barrido de puertas. Las rutas de traslado y montaje verifican embalaje, anclajes y acceso de servicio; ningún gabinete invade el acceso o la zona climática. La circulación útil no se acredita sólo con el área total. Cada climatizador dispone de envolvente 1,20$\times$0,70$\times$2,20 m. La circulación útil y el frente de servicio serán al menos 1,20 m; se respetan también las holguras adicionales de fabricantes. Las rutas A/B tienen canalizaciones y acceso propios. El stock local se mantiene apagado en recinto de respaldo separado, y su bancada portátil se energiza en el área externa de trabajo, con salida eléctrica y calor fuera del volumen climatizado \parencites{p7-eaton-93t60}{p6-mext-databook}{p12-apc-ar3104}.

Para acotar la envolvente/ocupación a 0,500 kW de aporte sensible se prescribe transmitancia global de cerramientos $U\leq0,120$ W/(m$^2$ K), incluidos puentes, y protección solar. Con $A=2(70)+2(10+7)(3)=242$ m$^2$, salto térmico de 11 K, fuga equivalente 15 m$^3$/h, densidad conservadora 1,2 kg/m$^3$ y calor específico 1,020 kJ/(kg K), la transmisión y fuga son $0,31944+0,05610=0,37554$ kW. Reservando 0,020 kW solares y 0,100 kW de ocupación técnica, el total es 0,49554 kW. Son límites propios de nueva obra, sujetos a cálculo térmico, sellado y recepción; no caracterizan el inmueble actual. La carga latente y los estados de baterías/clima se desarrollan por separado en 4.2.1. Los 15 m$^3$/h describen fugas residuales, sin sumar la renovación del banco. Las dos zonas de baterías tendrán extracción y reposición propias, con separación de aire respecto de cómputo/comunicaciones; el tratamiento de esa reposición se selecciona desde la memoria por estados. El módulo de 70 m$^2$ conserva carácter de implantación proyectada y se ajusta a la configuración ambiental elegida, manteniendo las holguras y el dimensionamiento proporcional.

Cada zona exterior de generación conserva una plataforma exclusiva de al menos 6,0$\times$4,0 m para el grupo, sin situar allí tanque ni bancada Avtron. El depósito de obra nueva tiene 2,00$\times$1,00$\times$1,00 m interiores y 2.000 L nominales; 200 L de expansión y hasta 300 L no extraíbles conservan al menos 1.500 L útiles por camino. El cubeto ofrece 2.200 L netos como mínimo: dimensión interior de proyecto 3,00$\times$2,00 m y altura útil de líquido de 0,60 m, con pared de 0,70 m. A esa altura, 3,60 m$^3$ brutos menos 1,20 m$^3$ de desplazamiento del tanque y hasta 0,20 m$^3$ de soportes dejan 2,20 m$^3$ netos. Se verifican desplazamientos, impermeabilidad, resistencia, drenaje controlado y normativa de combustible antes de ejecutar; el volumen de proyecto no constituye estanque certificado. La transferencia exterior por camino es de al menos 320 A y cuatro polos según esquema aprobado, sin cortar PE ni confundirla con el MCCB del grupo \parencites{p7-eaton-93t-guia}{p7-fg-p200-ficha}.

\subsubsection{Blindaje y resistencia de la envolvente no estructural}
\label{sec:recinto-blindaje-rt0602}
Synaptix fija la especificación BPR-4 1.0 para los cuatro muros no estructurales del módulo, incluidos encuentros con solera, cubierta, accesos y penetraciones. La amenaza de diseño es intrusión para robo o sabotaje mediante palanca, martillo, cincel y taladro portátil. Se prescribe una piel continua interior de acero estructural S355JR de 6 mm nominales, con espesor recibido no menor de 6 mm y certificado de material EN 10204 tipo 3.1; bastidor de perfiles tubulares 80$\times$80$\times$4 mm a paso máximo de 0,60 m y fijación a estructura resistente, sin transferir la resistencia a placas livianas o a los racks. Soldaduras continuas en encuentros y juntas de paneles, protegidas contra corrosión e inspeccionadas, conservan la continuidad; los elementos desmontables se fijan desde el interior y no dejan tornillos atacables desde fuera. El constructor calculará perfiles, soldaduras, anclajes y acción sísmica sobre la estructura nueva. Estas dimensiones constituyen una especificación propia de obra, sin atribuir un grado de resistencia a un producto no ensayado.

Se exige resistencia de recepción de al menos diez minutos de ataque efectivo sin apertura pasante que permita acceso corporal ni liberación del cierre desde fuera, sobre paño completo, junta, anclaje, encuentro superior/inferior y penetración representativos. El protocolo acordado antes de fabricación usa el método y herramientas de ataque manual de EN 1630 y documenta tiempos efectivos, herramienta, cara de ataque, puntos ensayados, criterio de paso y resultado. EN 1627 clasifica puertas, ventanas, rejas y cerramientos enumerados en esa norma; no se declara una clase RC4 universal para este muro por usar una chapa de 6 mm. Las puertas exteriores se suministrarán como conjuntos RC4 conforme a EN 1627:2021, con clasificación y ensayos de la misma variante, dimensiones, herrajes y montaje, además del ensayo de su unión al muro. El laboratorio independiente comprobará el conjunto BPR-4 antes de autorizar fabricación en serie y la recepción verificará que coincide con la muestra. La oferta no afirma un certificado de blindaje obtenido \parencite{lote10-bs-en1627}.

La continuidad no se pierde por el aire ni el cableado: pasamuros de dimensión mínima para cada servicio, manguitos de acero de 6 mm, fijación interior y sello; rejillas de ventilación con laberinto y barrera metálica, sin acceso recto a mecanismos o equipos. Cada paso conserva su resistencia de intrusión y el cierre de fuego que corresponda mediante un detalle ensayado; no se reemplaza un sistema cortafuego por una soldadura. Las salidas abren desde dentro sin credencial, llave ni energía, con luz útil de 1,20 m; la protección contra intrusión no retrasa evacuación. Se coordinarán fuerza de apertura, herraje de emergencia, control de acceso y compartimentación antes de compra.

La superficie bruta de chapa perimetral, sin descontar huecos ni añadir bastidor, es $2(10+7)(3)=102$ m$^2$: $102(0,006)(7.850)=4.804,2$ kg de cota propia de material. Esta carga se incorpora a la estructura y cimentación; no se considera dentro de los 25 kN/m$^2$ previstos para los bancos ni se atribuye al edificio actual. El aislamiento exterior continuo de lana mineral de 350 mm, con conductividad declarada $\leq0,034$ W/(m K), tiene referencia plana $0,034/0,350=0,0971$ W/(m$^2$ K). El proyecto térmico deberá demostrar el límite global anterior $U\leq0,120$, incluidos bastidor, puertas, cubierta, piso y puentes; el cociente aislante no lo acredita. Las dimensiones de 10$\times$7 m son interiores libres después del blindaje y aislamiento. No se recorta ese espacio para acomodar el muro. La resistencia frente a intrusión se recibe separadamente de fuego, estanqueidad, corrosión y aislamiento.

\subsubsection{Implantación coordinada de zonas y recorridos}
\label{sec:recinto-implantacion-coordinada}
La Figura~\ref{fig:recinto-exterior-coordinado} y las coordenadas siguientes fijan una reserva de terreno nueva de 30$\times$24 m, con origen en su esquina suroeste, norte hacia arriba y cotas en metros. Son dimensiones de proyecto a materializar, no un levantamiento de la marina. El interior del módulo ocupa $x=10$--20, $y=8$--15; los espesores de muro se añaden hacia fuera y las separaciones se miden entre caras exteriores finales. El plano interno usa su propio origen en la esquina suroeste del interior; por tanto $x_{sitio}=x_{interno}+10$ y $y_{sitio}=y_{interno}+8$.

\begin{longtable}{L{\dimexpr.23\linewidth-2\tabcolsep\relax}L{\dimexpr.30\linewidth-2\tabcolsep\relax}L{\dimexpr.47\linewidth-2\tabcolsep\relax}}
\caption{Zonas de obra nueva y reservas de circulación}\label{tab:recinto-coordenadas}\\\hline
\EncabezadoTabla{Zona} & \EncabezadoTabla{Coordenadas de proyecto} & \EncabezadoTabla{Holgura, acceso y separación}\\\hline
\endfirsthead
\hline\EncabezadoTabla{Zona} & \EncabezadoTabla{Coordenadas de proyecto} & \EncabezadoTabla{Holgura, acceso y separación}\\\hline\endhead
Generación A/B & A: $[0,6]\times[18,22]$; B: $[24,30]\times[18,22]$ & Dos plataformas 6$\times$4 m; envolvente del grupo y al menos 1,20 m de mantenimiento dentro de ellas deben ser compatibles. Escape elevado y dirigido fuera de accesos, admisiones y descarga de bancos; no se recibe si recircula gases.\\\hline
Combustible A/B & A: $[0,4]\times[10,13]$; B: $[26,30]\times[10,13]$ & Tanque por camino de 2.000 L nominales y contención de 2.200 L netos. Reserva de 4$\times$3 m incluye acceso de reposición; selección de tanque/cubeto verifica volumen neto y distancias de incendio aplicables. No se sitúan en la plataforma del grupo.\\\hline
Unidades exteriores clima & A: $[12,14]\times[19,21]$; B: $[16,18]\times[19,21]$ & Admisión/descarga y frente de 1,20 m libres, sin apuntar a puerta norte de bancos. Tuberías independientes hacia las dos unidades interiores; longitud/carga R32 se selecciona para este recorrido y no se presume el punto de catálogo de 5 m.\\\hline
Trabajo y enrolamiento & $[12,15]\times[1,3]$ & Zona cerrada y controlada de 3$\times$2 m, con bancada fuera de TI, puesto/lector de enrolamiento y paso propio a acceso. No habilita por sí sola RT-06.21.\\\hline
Respaldo local & $[17;19]\times[1;2{,}5]$ & Zona cerrada de 2$\times$1,5 m, stock apagado, separada del combustible y sin equiparar este stock a copias de datos.\\\hline
Escape y servicio & Frente sur: $y=5$--8; norte: $y=15$--18; laterales: $x=7$--10 y 20--23 & Corredores exteriores continuos libres de al menos 1,50 m después de muros, sin depósitos, bandejas bajas, cerramientos bloqueados o invasión de hojas. Frente TI sale al sur; ambos bancos salen al norte y las UPS tienen salida lateral.\\\hline
\end{longtable}
\FuenteTabla{Diseño de Synaptix: zonas y separaciones mínimas propias, sometidas al cálculo civil, incendio, ventilación, combustible y requisitos de montaje de las variantes suministradas.}

El grupo P200-6 de la ficha consultada mide 2,510$\times$1,010$\times$1,640 m y pesa 1.593 kg húmedo, sin atribuir esa huella a una cabina acústica opcional: con 1,20 m libres alrededor requiere 4,910$\times$3,410 m y cabe geométricamente en la plataforma de 6$\times$4 m. El montaje exterior necesita seleccionar protección ambiental/acústica compatible y recibirla; la ficha no acredita una envolvente marina. El fabricante publica 328 m$^3$/min de aire de radiador, 125 Pa de restricción externa y gases de escape a 530\,$^{\circ}$C: el trazado de admisión/escape y sus pérdidas se calculará, sin asumir que la distancia dibujada disipa esa carga \parencite[pp.~1--3]{p7-fg-p200-ficha}.

Entre las reservas de combustible y el módulo hay 6 m nominales en planta antes de espesores; el grupo más próximo queda a 5 m de la esquina interior del módulo. Esas cotas no se presentan como distancias reglamentarias suficientes: el proyecto de seguridad valida combustible, ruido, evacuación, incendio y viento, y ampliará la reserva exterior si su norma o fabricante exige más. El abastecimiento recorre el perímetro exterior oeste/este hasta cada cubeto y no cruza las salidas ni el interior. A y B poseen ductos, pasamuros, tablero, transferencia y ruta propios por lados opuestos; una excavación o soporte único no puede cortar ambos.

Los bancos mantienen extracción y reposición separadas del aire TI. Se reserva conexión independiente por banco hacia su fachada norte, con descarga elevada y sin retorno a climatización; selección de ajuste real de cargador, calor, caudal, acondicionamiento y clasificación de ventilador continúa abierta en la memoria por estados. Ninguna flecha del plano acredita que el escenario de 300/900 m$^3$/h sea el definitivo. Se exige mantener paso de 1,30 m frente a cada cadena y de 1,20 m en UPS, racks y climatizadores, además de los mínimos de fabricante. Las UPS se desplazan hacia su borde exterior y abren al lateral; los racks ocupan huellas separadas de 0,80$\times$1,00 y 0,60$\times$0,80 m, con frente de 1,20 m dentro de TI. El diseño ejecutivo verificará también acceso posterior de racks, traslado de bloques, grúa y carga de mantenimiento. RT-06.03 conserva desarrollo pendiente hasta seleccionar el tratamiento de bancos y verificar envolventes/holguras de variantes; se entrega distribución de oferta, sin afirmar plano recibido.

\begin{figure}[htbp]
\centering
\setstretch{1}
\begin{tikzpicture}[x=.42cm,y=.42cm,
 zona/.style={draw=SynNavyLight,fill=SynSurface,thick},
 texto/.style={font=\fontsize{8.5}{10}\selectfont,align=center},
 ruta/.style={draw=SynBlue,thick,dashed},
 cota/.style={{Latex}-{Latex},thin}]
\draw[draw=SynNavyLight,dashed] (-2,-2) rectangle (32,37);
\draw[zona] (9.5,7.5) rectangle (20.5,15.5);
\node[texto,text width=39mm] at (15,11.5) {Módulo TI/energía\\10$\times$7 m interiores\\11$\times$8 m exteriores\\Plano interno ampliado};
\draw[zona] (0,18) rectangle (6,22);
\node[texto,text width=23mm] at (3,20) {Grupo A\\6$\times$4 m};
\draw[zona] (24,18) rectangle (30,22);
\node[texto,text width=23mm] at (27,20) {Grupo B\\6$\times$4 m};
\draw[zona] (0,10) rectangle (4,13);
\node[texto,text width=17mm] at (2,11.5) {Tanque A\\y cubeto};
\draw[zona] (26,10) rectangle (30,13);
\node[texto,text width=17mm] at (28,11.5) {Tanque B\\y cubeto};
\draw[zona] (12,19) rectangle (14,21);
\node[texto,anchor=south] at (13,21.1) {Clima A};
\draw[zona] (16,19) rectangle (18,21);
\node[texto,anchor=south] at (17,21.1) {Clima B};
\draw[zona] (12,1) rectangle (15,3);
\node[texto,anchor=north] at (13.5,.9) {Trabajo\\3$\times$2 m};
\draw[zona] (17,1) rectangle (19,2.5);
\node[texto,anchor=north] at (18,.9) {Respaldo\\2$\times$1,5 m};
\draw[-{Latex},thick] (15,8)--(15,5);
\draw[-{Latex},thick] (11.8,15)--(11.8,17);
\draw[-{Latex},thick] (18.2,15)--(18.2,17);
\draw[-{Latex},thick] (10,9.1)--(7,9.1);
\draw[-{Latex},thick] (20,9.1)--(23,9.1);
\node[texto,text width=39mm] at (15,4.7) {Escape sur / pasos laterales\\libres $\geq1,50$ m};
\node[texto,text width=39mm] at (15,17.3) {Escape bancos al norte\\extracción propia elevada};
\draw[ruta] (3,18)--(7,18)--(7,10)--(10,10);
\draw[ruta] (27,18)--(23,18)--(23,10)--(20,10);
\node[texto,rotate=90,fill=white] at (7,13) {Ruta A};
\node[texto,rotate=90,fill=white] at (23,13) {Ruta B};
\draw[cota] (-2,36)--(32,36);
\node[texto,fill=white] at (15,36) {Reserva de proyecto: 34 m};
\draw[cota] (-2.7,-2)--(-2.7,35);
\node[texto,rotate=90,fill=white] at (-2.7,16.5) {37 m};
% Huellas de variante y reservas de obra coordinadas ENV-19/PORT-21.
\draw[zona] (0,0) rectangle (4,7);
\draw[zona] (26,0) rectangle (30,7);
\node[texto,text width=16mm] at (2,3.5) {Bancada A\\4$\times$7 m};
\node[texto,text width=16mm] at (28,3.5) {Bancada B\\4$\times$7 m};
\draw[zona] (6.5,24) rectangle (11.5,28);
\draw[zona] (18.5,24) rectangle (23.5,28);
\draw[fill=white,draw=SynNavyLight] (7.88,25.42) rectangle (10.12,26.57);
\draw[fill=white,draw=SynNavyLight] (19.88,25.42) rectangle (22.12,26.57);
\node[texto,text width=20mm] at (9,27.3) {Frío A\\30RB065-SS};
\node[texto,text width=20mm] at (21,27.3) {Frío B\\30RB065-SS};
\draw[zona] (7.5,20) rectangle (10.5,22);
\draw[zona] (6,31) rectangle (9,33);
\draw[zona] (19.5,20) rectangle (22.5,22);
\draw[zona] (21,31) rectangle (24,33);
\node[texto] at (9,21) {UTA A1};
\node[texto] at (7.5,32) {UTA A2};
\node[texto] at (21,21) {UTA B1};
\node[texto] at (22.5,32) {UTA B2};
\draw[zona] (7.5,4) rectangle (9.5,6);
\draw[zona] (20.5,4) rectangle (22.5,6);
\node[texto,text width=12mm] at (8.5,5) {Agua A};
\node[texto,text width=12mm] at (21.5,5) {Agua B};
\draw[zona] (11,31) rectangle (14,33);
\draw[zona] (16,31) rectangle (19,33);
\node[texto,text width=12mm] at (12.5,32) {Cilindros\\TI/bancos};
\node[texto,text width=12mm] at (17.5,32) {Reserva\\pequeños};
\draw[draw=SynTeal,dash dot] (9,24)--(8,23)--(8,16)--(10,12);
\draw[draw=SynTeal,dash dot] (21,24)--(22,23)--(22,16)--(20,12);
\node[texto,fill=white,inner sep=.5mm] at (8,16.4) {Agua A};
\node[texto,fill=white,inner sep=.5mm] at (22,16.4) {Agua B};
\draw[draw=SynNavyLight,fill=white] (10.2,11.2) rectangle (11.4,12);
\draw[draw=SynNavyLight,fill=white] (10.2,13.2) rectangle (11.4,14);
\draw[draw=SynNavyLight,fill=white] (18.6,11.2) rectangle (19.8,12);
\draw[draw=SynNavyLight,fill=white] (18.6,13.2) rectangle (19.8,14);
\node[texto,anchor=east] at (10.1,13.6) {EX A};
\node[texto,anchor=west] at (19.9,13.6) {EX B};
% Bastidores en cubierta, no sobre los pasos de salida norte.
\draw[zona] (7.75,17) rectangle (10.25,18);
\node[texto,font=\fontsize{7}{8}\selectfont] at (9.0,17.5) {IMP A1};
\draw[zona] (19.75,17) rectangle (22.25,18);
\node[texto,font=\fontsize{7}{8}\selectfont] at (21.0,17.5) {IMP B1};
\draw[zona] (6.25,34.5) rectangle (8.75,35.5);
\node[texto,font=\fontsize{7}{8}\selectfont] at (7.5,35.0) {IMP A2};
\draw[zona] (21.25,34.5) rectangle (23.75,35.5);
\node[texto,font=\fontsize{7}{8}\selectfont] at (22.5,35.0) {IMP B2};
\draw[draw=SynBlue,densely dotted,line width=.5pt] (3.0,17.25)--(-0.75,17.25)--(-0.75,20.0); % PORT28 grupoA
\draw[draw=SynBlue,densely dotted,line width=.5pt] (27.0,17.25)--(30.75,17.25)--(30.75,20.0); % PORT28 grupoB
\draw[draw=SynBlue,densely dotted,line width=.5pt] (-0.75,3.5)--(-0.75,12.25); % PORT28 tanqueA
\draw[draw=SynBlue,densely dotted,line width=.5pt] (30.75,3.5)--(30.75,12.25); % PORT28 tanqueB
\draw[draw=SynBlue,densely dotted,line width=.5pt] (-0.75,3.5)--(-0.75,-0.75)--(2.0,-0.75); % PORT28 bancadaA
\draw[draw=SynBlue,densely dotted,line width=.5pt] (25.25,3.5)--(25.25,-0.75)--(28.0,-0.75); % PORT28 bancadaB
\draw[draw=SynBlue,densely dotted,line width=.5pt] (11.25,20.0)--(11.25,18.25)--(14.75,18.25)--(14.75,20.0); % PORT28 frioTI_A
\draw[draw=SynBlue,densely dotted,line width=.5pt] (18.75,20.0)--(18.75,18.25)--(15.25,18.25)--(15.25,20.0); % PORT28 frioTI_B
\draw[draw=SynBlue,densely dotted,line width=.5pt] (3.0,22.75)--(3.0,23.25)--(8.5,23.25); % PORT28 frioA
\draw[draw=SynBlue,densely dotted,line width=.5pt] (18.75,20.0)--(18.75,23.25)--(21.5,23.25); % PORT28 frioB
\draw[draw=SynBlue,densely dotted,line width=.5pt] (11.25,20.0)--(11.25,19.25)--(6.75,19.25)--(6.75,20.25); % PORT28 UTA_A1
\draw[draw=SynBlue,densely dotted,line width=.5pt] (18.75,20.0)--(18.75,19.25)--(23.25,19.25)--(23.25,20.25); % PORT28 UTA_B1
\draw[draw=SynBlue,densely dotted,line width=.5pt] (7.5,36.25)--(5.5,36.25)--(5.5,33.75)--(7.5,33.75); % PORT28 UTA_A2
\draw[draw=SynBlue,densely dotted,line width=.5pt] (22.5,36.25)--(20.5,36.25)--(20.5,33.75)--(22.5,33.75); % PORT28 UTA_B2
\draw[draw=SynBlue,densely dotted,line width=.5pt] (4.75,3.5)--(4.75,6.75)--(10.25,6.75); % PORT28 aguaA
\draw[draw=SynBlue,densely dotted,line width=.5pt] (16.0,3.75)--(16.0,6.75)--(20.75,6.75); % PORT28 aguaB
\draw[draw=SynBlue,densely dotted,line width=.5pt] (7.5,36.25)--(9.5,36.25)--(9.5,33.75)--(10.75,33.75); % PORT28 cilTI
\draw[draw=SynBlue,densely dotted,line width=.5pt] (15.0,32.0)--(15.0,33.75)--(19.0,33.75); % PORT28 cilPeq
\draw[draw=SynBlue,densely dotted,line width=.5pt] (3.0,17.25)--(3.0,16.25)--(8.5,16.25); % PORT28 imp_A1
\draw[draw=SynBlue,densely dotted,line width=.5pt] (18.75,20.0)--(18.75,16.25)--(21.5,16.25); % PORT28 imp_B1
\draw[draw=SynBlue,densely dotted,line width=.5pt] (7.5,36.25)--(5.5,36.25)--(5.5,33.75)--(7.5,33.75); % PORT28 imp_A2
\draw[draw=SynBlue,densely dotted,line width=.5pt] (22.5,36.25)--(20.5,36.25)--(20.5,33.75)--(22.5,33.75); % PORT28 imp_B2
\draw[draw=SynBlue,densely dotted,line width=.5pt] (16.0,3.75)--(12.0,3.75)--(12.0,1.0); % PORT28 trabajo
\draw[draw=SynBlue,densely dotted,line width=.5pt] (16.0,3.75)--(16.0,1.0)--(19.0,1.0); % PORT28 stock
\foreach \px/\py/\pid in {3/17.25/10,27/17.25/11,4.75/11.5/12,25.25/11.5/13,4.75/3.5/14,25.25/3.5/15,12.25/26/16,17.75/26/17,16/3.75/18,3/22.75/19,27/22.75/20,-0.75/3.5/21,30.75/3.5/22,7.5/30.25/23,22.5/30.25/24,15/32/25,11.25/20/26,18.75/20/27,7.5/36.25/28,22.5/36.25/29}{
 \node[circle,draw=SynNavyLight,fill=white,inner sep=.7mm,font=\fontsize{7}{8}\selectfont] at (\px,\py) {P\pid};
}
\node[texto,text width=64mm] at (15,36.3) {P10--P29: portátiles exteriores; P01--P09: plano interno.\\UTA y agua: reservas de obra, conjuntos a medida pendientes.};
\draw[-{Latex},thick] (29,2)--(29,5);
\node[texto] at (29,5.6) {N};
\end{tikzpicture}
\caption{Reserva exterior y separación de zonas del módulo técnico nuevo}
\label{fig:recinto-exterior-coordinado}
\FuenteFigura{Elaboración de Synaptix: coordenadas de proyecto, sin levantamiento ejecutado. Rectángulos: reservas de equipos; flechas: escape; discontinuas: familias de ruta A/B, con datos y electricidad en canalizaciones separadas. Las holguras finales se miden desde caras exteriores, y la selección ambiental, combustible e incendio puede ampliar estas reservas. PORT-24 amplía la reserva a 34 por 39 m, desde $x=-2$ hasta 32 y $y=-2$ hasta 37; las huellas interiores de frío son de catálogo, las UTA/agua y recipientes pequeños son reservas de proyecto. EX A/B son bastidores en cubierta, fuera de salidas. Verde: familia hidráulica segregada; el trazado ejecutivo verifica drenaje, apoyos, cruces, gas y distancias, sin considerar estas líneas tuberías recibidas. La reserva de muro de 0,50 m por lado concilia blindaje, bastidor y aislamiento; la reproducción de rutas conserva máximo 8,875 m. PORT-24 representa las estaciones y rutas críticas con ejes a 0,75 m de las reservas; las reservas de máquina incluyen holguras y no se cruzan para recuperar un portátil.}
\end{figure}


\subsubsection{Canalizaciones, mapa de enlaces y certificación individual}
\label{sec:recinto-enlaces-certificados}
La especificación CER-4 1.0 fija solera continua sin piso elevado: el peso de los bancos se apoya directamente en estructura calculada y los servicios transcurren por bandejas aéreas accesibles, a 2,40 m o más, sin invadir puertas ni el frente de mantenimiento. Esta adaptación de la tipología secundaria se declara expresamente; no supone dispensar RT-06.04. Si las bases o la recepción exigen piso técnico elevado, su selección deberá acreditar cargas uniformes/concentradas, sísmicas, reacción al fuego y paneles desmontables, sin colocar bancos sobre un sistema sin cálculo. Ese detalle y su aceptación siguen pendientes; no se declara un piso técnico inexistente como instalado.

Electricidad aplica Decreto 8 y pliegos SEC RIC 04 (conductores/canalizaciones), RIC 06 (tierra), RIC 11, sección 15 (muelles), RIC 18 y RIC 19 (proyecto/puesta en servicio). RIC 11, 15.3 exige conexiones por encima del muelle, canalizaciones para lugares mojados, soportes estructurales y alivio de esfuerzos; en muelle flotante se diseña conexión al menos 300 mm sobre plataforma, y en fijo al menos 300 mm sobre plataforma y 600 mm sobre la mayor marea. Los enchufes a embarcaciones conservan diferencial de sensibilidad no mayor de 30 mA conforme a 15.3.5; esto no se extrapola sin cálculo a alimentadores UPS. Datos aplica ISO/IEC 11801-1 y -3, ISO/IEC 14763-2 (instalación), IEC 61935-1 (ensayo cobre), ISO/IEC 14763-3 (fibra) y ANSI/TIA-606-D (administración). IEC 60529 verifica envolventes; IEC 60068-2-52 verifica aptitud frente a niebla salina del conjunto seleccionado. Las ediciones y revisiones vigentes se consignarán en el proyecto aprobado; la certificación de datos no sustituye la declaración eléctrica \parencites{lote10-sec-ric04}[secc.~15, p.~58]{lote10-sec-ric11}{lote10-sec-ric}.

Se prescriben tres familias de ruta: datos A/B en bandejas metálicas con tapa de 200$\times$60 mm independientes por fachada oeste/este, electricidad A/B en bandejas de 300$\times$100 mm distintas y buses de instrumentación en canaleta propia de 100$\times$60 mm. Dimensiones son reservas de obra, verificadas por ocupación, capacidad térmica, radios y carga con la BOM. Entre energía y datos se fija separación de proyecto de 300 mm y cruce a 90 grados; si el estudio EMC o RIC requiere más, prevalece la mayor distancia. No comparten tabique, soporte colapsable, registro ni penetración A/B. Ductos enterrados de datos HDPE de 63 mm por camino, con registro estanco y drenaje controlado, se calculan según terreno y ocupación; las rutas eléctricas enterradas conservan ductos/registro propios y distancias RIC. Acero/soportes expuestos son AISI 316L o sistema anticorrosivo con aptitud marina comprobada; se equipotencializan las partes conductoras y se aíslan metales incompatibles.

La fibra seleccionada es OS2 monomodo, doce fibras por cable troncal, cubierta exterior resistente a UV/agua y tramo móvil de fabricación apta para flexión repetida. El cobre fijo es Cat 6A F/UTP de cobre sólido, libre de halógenos dentro de sala, con continuidad de pantalla y tierra diseñadas; el enlace permanente no supera 90 m y el canal completo 100 m. No se usa cable sólido fijo en la articulación tierra-pantalán. Ese tramo utiliza cable dinámico y funda/guiado flotante seleccionados por radio dinámico, torsión, carga de tracción, temperatura, salinidad y número de ciclos declarado; empalmes y conectores permanecen en cajas protegidas por encima de plataforma. Ninguna ficha de ``uso exterior'' acredita flexión permanente.

\begin{paginaHorizontalSynaptix}
\begingroup
\setlength{\tabcolsep}{3pt}
\begin{longtable}{L{\dimexpr.17\linewidth-2\tabcolsep\relax}L{\dimexpr.29\linewidth-2\tabcolsep\relax}L{\dimexpr.15\linewidth-2\tabcolsep\relax}L{\dimexpr.39\linewidth-2\tabcolsep\relax}}
\caption{Mapa CER-4 de enlaces de diseño, cantidades y aceptación}\label{tab:cer4-enlaces}\\\hline
\EncabezadoTabla{ID/familia} & \EncabezadoTabla{Extremos y medio} & \EncabezadoTabla{Cantidad} & \EncabezadoTabla{Longitud, ruta y documento individual}\\\hline\endfirsthead
\hline\EncabezadoTabla{ID/familia} & \EncabezadoTabla{Extremos y medio} & \EncabezadoTabla{Cantidad} & \EncabezadoTabla{Longitud, ruta y documento individual}\\\hline\endhead
TR-P01..05-A/B; TR-V01-A/B & Core A/B--switch A/B de cinco pantalanes y varadero. OS2, dúplex por enlace; doce fibras/cable & 12 cables/12 enlaces; 24 fibras activas, 120 reserva & Cota de proyecto 1.000 m/cable, 12 km máximos asignados antes de medir. A y B separados; dos por gabinete. Certificado óptico por cada par y registro de las fibras de reserva, sin certificarlas como enlace activo.\\\hline
CN-A/B-01..06 & Terminación de troncal--switch del mismo gabinete, conversor óptico si lo exige puerto seleccionado & 12 conversores dentro de la cota de energía existente & Una terminación remota por troncal; extremo core mediante SFP compatible. Sin duplicar un conversor cuando el switch incorpora SFP; BOM, doce pares de ópticas y sus reservas compatibles por cerrar.\\\hline
LC-S01..02-A/B & Dos nodos--dos core; cobre Cat 6A & 4 enlaces & $\leq15$ m/enlace, rutas A/B. Red de servicio/réplica/gestión segregada por VLAN; no se atribuye un puerto adicional inexistente. Cada NIC/puerto y HA se valida.\\\hline
LC-F01..02-A/B; LC-HA-01 & Dos NGFW--core y enlace HA dedicado; Cat 6A & 4 +1 enlaces & $\leq15$ m; dos core por NGFW; HA separado. Firmware, puerto y certificación cobre individual.\\\hline
LC-ST-01/02 & Stack core--core; cable de stack propio fabricante & 2 enlaces físicos & $\leq3$ m. Referencia/longitud y compatibilidad del kit; ensayo fabricante, no certificado Cat 6A inventado.\\\hline
LC-WAN-A/B-01..02 & CPE de cada operador--puertos WAN propios en ambos NGFW; Cat 6A o entrega óptica documentada & 4 enlaces de entrega & $\leq30$ m; el operador define interfaz y permite distribución VLAN independiente. La ruta exterior WAN se documenta en su contrato, no en los 12 km de troncal local.\\\hline
PT-01..16 & Core/patch panel--puestos fijos & 16 enlaces Cat 6A & $\leq90$ m permanentes; $\leq100$ m canal. Reserva propia 16$\times$90=1.440 m, sin contar latiguillos dentro del permanente. Sin segundo cable ficticio para un PC de NIC única.\\\hline
AP-P01..05-A/B; AP-V01-A/B & Switch de gabinete--dos AP por sitio, cobre PoE & 12 enlaces iniciales & $\leq15$ m fijos/enlace. Si AP final usa dos puertos, segundo enlace se añade a BOM y cálculo, sin asumirlo por la pareja de switches. Certificación cobre y caída de tensión PoE.\\\hline
CT-P01..05-01/02-A/B & Switch A/B--dos concentradores por pantalán & 20 enlaces propuestos & $\leq15$ m. Diez concentradores con dos interfaces sólo si la variante lo admite; acceso alternativo y aislamiento por validar, no duplicar concentradores en inventario.\\\hline
ACC-01..07; COM-01; MET-01 & Controladores de siete accesos, interfaz segura combustible y estación meteo--distribución autorizada & 9 enlaces previstos & $\leq90$ m cobre fijo por enlace o sustitución óptica si excede. Combustible: ruta sólo hasta interfaz en oficina segura; cualquier paso a zona clasificada requiere barrera y certificado de esa variante. Protocolos y puertos pendientes.\\\hline
CCTV-001..038 /039..060 & Cámara--plataforma CCTV/puerto de su segmento & 38 existentes; reserva hasta 60 & No se certifican por presunción: inspección individual de medio, recorrido, alimentación y ensayo. Longitudes/adaptaciones por levantar; canales que excedan límites se reemplazan con fibra y cota propia antes de compra. NVR--core: enlace adicional CCTV-NVR-01.\\\hline
BUS-Ppp-Ccc-W/E-bb; MW/ME-Ppp-nnn & Concentrador--buses agua/energía--medidor de cada amarra, RS-485 propuesto & 40 buses iniciales/60 crecimiento; 640/760 derivaciones de medidor & Cinco pantalanes: 64/76 amarras cada uno, dos concentradores; 32/38 medidores por servicio/concentrador. Máximo 16 medidores/bus: dos/tres buses por servicio. $\leq200$ m por bus, $\leq0,30$ m derivación; cable industrial par 120 ohm, terminación en ambos extremos y un solo punto de referencia según transceptores aislados. Se prueba cada bus y cada lectura individual; no se presenta como enlace Ethernet certificado.\\\hline
PIL-01..040; AUX-xx & Sensores piloto y seguridad/ambiente de sala, incendio, PDU/BMC y radios según interfaz de fabricante & 40 puntos piloto; auxiliares por BOM & Medio, puertos, cantidades y recorrido final aún abiertos. Ninguna suma de cables de la tabla cubre interfaces no seleccionadas; se integra en CER-4 antes de aceptar RT-06.04.\\\hline
\end{longtable}
\FuenteTabla{Diseño de Synaptix desde T-11: cantidades de anfitriones conservadas; las longitudes son cotas propuestas y las familias pendientes no se contabilizan como inventario definitivo.}
\endgroup
\end{paginaHorizontalSynaptix}

Para los buses, los 320 medidores de agua y 320 de energía iniciales se distribuyen como 64 por servicio en cada pantalán; crecimiento: 380 por servicio, 76 por pantalán. Cada concentrador requiere cuatro canales RS-485 aislados al inicio y seis en crecimiento; una variante con menos puertos exige multiplexores/aisladores y actualizar T-11/energía antes de compra. Las 640/760 derivaciones no se suman como troncales Ethernet ni convierten la propuesta RS-485 en un atributo de catálogo de un medidor sin seleccionar. La interfaz del totalizador y aptitud marina siguen pendientes. Los 12 conversores se mantienen dentro de la hoja de cargas, no se agregan a ella nuevamente.

La familia de ID identifica cada unidad, no sólo el grupo: por ejemplo TR-P03-B, CT-P03-02-A y ME-P03-045. El registro CER-4 conserva etiqueta legible en ambos extremos, puerto/patch panel, fibra/par, gabinete, camino, ruta, fecha, longitud diseñada y medida, marca/referencia/lote, norma/perfil, calibración del equipo, técnico, resultado y archivo original firmado/hash. Se entrega inventario abierto CSV y planos PDF/CAD; una etiqueta renumerada conserva historial. El enlace sin certificado o con fallo queda fuera de aceptación; no basta un porcentaje de muestras.

Cobre se certifica al 100\,\% como enlace permanente Cat 6A con wiremap, longitud, pérdida de inserción, NEXT/PSNEXT, retorno y demás parámetros del perfil de IEC 61935-1; se conserva prueba del canal y PoE cuando corresponda. Fibra se recibe con inspección/limpieza de conectores, polaridad, continuidad, longitud, pérdida óptica a 1.310/1.550 nm y presupuesto calculado desde transmisor/receptor, conectores, empalmes y margen; OTDR bidireccional complementa, sin sustituir OLTS. Todo resultado enlaza exactamente con el ID y sus dos extremos. Buses industriales conservan polaridad/continuidad, aislamiento adecuado a equipo, terminación, estabilidad y lectura/ACK de cada medidor bajo carga; cables de stack y enlaces de operador incluyen sus propios informes, sin aplicarles un certificado Cat 6A ajeno.

En cada paso tierra-pantalán se levantan elevaciones extremas, desplazamiento horizontal y oleaje antes de elegir el conjunto dinámico. Su longitud de diseño se calcula desde la máxima distancia entre anclajes, más reserva geométrica que conserve radio dinámico y alivio de tensión en todo el sobre; no se fija un porcentaje universal de holgura. Se conserva trazabilidad de ficha de cable dinámico y guiado, ciclos/ensayo de flexión y exposición salina, carga admisible y montaje real. La recepción recorre marea completa con oleaje de diseño, controla ausencia de tracción/estrangulamiento/roce e ingreso de agua y repite pérdida/errores después del ensayo. Inspección trimestral, revisión tras temporal y reserva compatible de tramo móvil por tipo, con umbral de sustitución según desgaste/ficha, se incluyen en el mantenimiento y costeo C07.

El CLIENTE adquiere hardware y ejecuta canalizaciones, obra y soportes; Synaptix especifica, coordina, integra y recibe técnicamente, y contrata servicios donde corresponda. CER-4 deberá completar longitudes exteriores levantadas, interfaces de todos los auxiliares/CCTV/piloto, cables dinámicos y BOM real de ópticas/medidores antes del cierre de RT-06.04. La certificación futura individual es una obligación concreta de recepción, sin afirmar ensayos ejecutados. El costeo económico identificará cada familia, metros de ducto/cable, soportes, cajas, etiquetas, certificados y reposición; la cota de 12 km no se ofrece como cantidad medida del sitio.


Synaptix entrega planos civil, estructural, unifilar, ventilación de baterías, canalizaciones y matriz de recepción. El CLIENTE ejecuta las obras asignadas; el instalador competente verifica soporte, tierra, protecciones, selectividad, enlaces y documentación chilena. Si el emplazamiento no cumple las condiciones de este módulo, se ajusta la implantación antes de compra, conservando separación y cálculos de la oferta.

GCI-4 individualiza la bancada Avtron 4100 de 60 kW/400 V/50 Hz y ALC por grupo, con una reserva compatible adicional. Su huella aproximada es 1,219$\times$0,685 m y su masa 227 kg; no cabe junto al grupo por presumir libre la plataforma de 6$\times$4 m. Se reserva otra plataforma de 4$\times$7 m por bancada: la dimensión 1,219+2(2,44)=6,099 m cabe en 7 m en el eje de admisión/descarga; el frente de mantenimiento permanece separado. Las ubicaciones exteriores definitivas, viento, rutas, protección marina y compatibilidad de motor/control a 50 Hz se incorporarán al plano ejecutivo; esta reserva adicional no acredita todavía el plano completo RT-06.03. ENV-16 e INC-4 requieren además huellas de tratamiento de bancos y cilindros, sin invadir escapes, tomas, combustible o evacuación \parencite{lote17-avtron-4100}.

\textbf{Sectorización INC-4 y auxiliares.} La circulación frontal de 5,0$\times$1,5 m se proyecta cerrada respecto de TI mediante el tabique representado en $y=1,5$ y una puerta de 1,20 m libre, de cierre supervisado y salida sin llave. La clasificación/resistencia del conjunto, holguras y alivio de presión se reciben con el proyecto de incendio; no se declara estanqueidad por dibujar el tabique. El gabinete PCI-4 y los cuatro lectores auxiliares se ubican fuera de bancos, en pared técnica de circulación, con proyección máxima propia de 0,25 m y sin invadir la franja central ni barridos; quedan al menos 1,25 m en ese tramo de paso. Gabinetes y bancos propios de protección deben ajustarse a esa cota o reubicarse antes de compra. Su cota térmica conjunta 0,300 kW normal/0,650 kW máxima se conserva dentro del módulo y de la frontera energética, en circulación; no se atribuye ese calor a la sala TI. ENV-16 debe incluir esa zona en su tratamiento N+1 por estado aún sin seleccionar.

\textbf{Coordinación de implantación ENV-19, EXT-21 y PORT-21.} La Figura~\ref{fig:recinto-exterior-coordinado} amplía la reserva exterior nueva a 30 por 33 m y conserva el módulo interior, grupos y combustible. Cada bancada dispone de su propia plataforma de 4 por 7 m en las esquinas sur; los chillers Carrier ocupan plataformas de 5 por 4 m en el norte, con huellas de catálogo representadas. Cuatro reservas UTA de 3 por 2 m alojan una envolvente propia de equipo de 2 por 1 m más acceso; no son dimensiones de catálogo. Dos áreas hidráulicas de 2 por 2 m conservan inercia y bombeo. Los bastidores EX de 1,2 por 0,8 m se sitúan elevados en cubierta sobre bancos, fuera del paso norte; las descargas elevadas, captaciones y protección del edificio se resuelven con gas y viento, conservando la separación propia de 3 m frente a tomas. No se libera el plano ejecutivo de RT-06.03 sin esas rutas, apoyos, variantes y holguras verificadas.

Las reservas de recipientes EXT-21 de 3 por 2 m se separan en el extremo norte: una aloja el gabinete de seis recipientes TI/bancos con su frente de servicio, y la otra conserva espacio para las tres zonas pequeñas aún sin variante. PORT-21 suma dieciocho portátiles activos y cuatro reservas locales: quince B441, un 398 y dos 322 activos; dos B441, un 398 y un 322 de reemplazo. Los puntos P01--P09 se muestran en el plano interno; P10--P18, en el exterior. La protección abarca también las plataformas ambientales, trabajo y stock; las reservas no se cuentan como posiciones activas. Se prescriben corredores perimetrales continuos de 1,20 m, sin cruzar cubetos ni huellas de máquinas, y recorrido efectivo máximo de 9 m con gabinete, puertas y barreras. Las áreas UTA/hidráulicas quedan dentro de las rutas de las estaciones ambientales o de energía próximas; sus riesgos reales y desvíos se cotejan con las variantes antes de ejecución. La recepción mide los recorridos y acredita potencial local/certificados conforme a DS 594 y DS 44. El diseño de portátiles no sustituye la detección ni extinción automática aún pendientes.
